\documentclass[letterpaper,10pt]{article} 
\usepackage{templateRJCP}
\newcommand\authoraffil[1]{\textsuperscript{#1}}
\renewcommand\refname{}
\usepackage{graphicx} % Required for inserting images
\usepackage[utf8x]{inputenc}
\begin{document}

%%%%%%%%%%%%%%%%%%%%%%%%%%%%%%%%%%%%%%%%%%%%%
\title{Template pour les Rencontres des Jeunes Chercheurs en Parole 2023}
\author{Prénom Nom,\authoraffil{1,*} Prénom Nom,\authoraffil{2} and Prénom Nom\authoraffil{2,3}}

\address{\authoraffil{1} affiliation 1\\
\authoraffil{2} affiliation 2\\
\authoraffil{3} affiliation 3}

\email{\authoraffil{*}adresse.mail@premierAuteur.fr} %% email address is required
%%%%%%%%%%%%%%%%%%%%%%%%%%%%%%%%%%%%%%%%%%%%%%%

\section*{Résumé}
Insérer ici le résumé de 300 mots maximum. Il peut être rédigé en français ou en anglais. Dans la section “mots-clés”, insérer au maximum 5 mots-clés, séparés par une virgule. Ajouter ensuite les thématiques auxquelles se rapportent le résumé proposé (au maximum 2, la plus importante en premier). Enfin, les références doivent apparaître dans la section “Références” \cite{TRAMERETAL_2023AcademicAuthorHow}.

\section*{Mots-clés}\label{sec:mots-cles}
parole, prosodie, rjcp (5 mots-clés maximum)

\section*{Thématiques}
* Acoustique de la parole
* Acquisition de la parole et du langage 
* Analyse, codage et compression de la parole
* Applications à composantes orales (dialogue, indexation, etc)
* Apprentissage d’une langue seconde
* Communication multimodale 
* Dialectologie 
* Évaluation, corpus et ressources 
* Langues en danger 
* Modèles de langage 
* Parole audio-visuelle
* Pathologies de la parole 
* Phonétique et phonologie 
* Phonétique clinique 
* Production / Perception de la parole
* Prosodie 
* Psycholinguistique 
* Reconnaissance et compréhension de la parole 
* Reconnaissance de la langue
* Reconnaissance du locuteur
* Signaux sociaux, sociophonétique 
* Synthèse de la parole
* Autre (préciser)

\section*{Langue utilisée pour la communication}
Précisez ici la langue de présentation du poster.


\section*{Références}
\bibliographystyle{plain} % We choose the "plain" reference style
\bibliography{references} % Entries are in the refs.bib file

\end{document}
